\documentclass[12pt, letterpaper]{article}
\usepackage{graphicx} % Required for inserting images
%\usepackage[T1]{fontenc}

\usepackage{amsfonts}
\usepackage{tabularx}
\usepackage{multirow}
\usepackage{booktabs}
\usepackage{float}

\usepackage{polski}
%\usepackage[utf8]{inputenc}

\title{Programy użytkowe \\ Lista 3}
\author{ }
\date{}

\begin{document}
\maketitle
\begin{enumerate}
    \item Umieść w dokumencie obrazek, wykorzystując któryś z przykładowych obrazów (np. \texttt{example-image-a}) zamiast własnoręcznie przesyłanego obrazu.
    \item Przetestuj różnicę w znaczeniu \verb|\inewidth| oraz \verb|\textwidth| przy ustalaniu szerokości obrazu (najlepiej w tym celu użyj przejdź do trybu dwukolumnowego).
    \item Użyj środowiska \verb|wrapfigure|, aby uzyskać efekt wykresu otoczonego tekstem (dodaj obrazek przesunięty do lewej i do prawej krawędzi).
    \item Umieść dwa wykresy obok siebie, do każdego z nich dodając osobny podpis (użyj \verb|minipage| lub \verb|subfigure|). 
    Poeksperymentuj z przerwami między wykresami (\verb|\hfill|, \verb|\hspace{}|)
    
    \item Stwórz tabelę, jak poniżej
    \begin{table}[H]
        \centering
          \caption{Podstawowe informacje o kilku funkcjach elementarnych}
        \begin{tabular}{c c c c}
        \toprule
        $f(x)$ & dziedzina & zbiór wartości & miejsce zerowe \\
        \midrule
        \midrule
        ${x^2}$ & $\mathbb{R}$ & $\mathbb{R}_+$ & $0$ \\
        $\cos(x)$ & $\mathbb{R}$ & $[-1,1]$ & $2k\pi + \frac{\pi}{2},\,k\in\mathbb{Z}$ \\
       $ \exp(x)$ & $\mathbb{R}$ & $\mathbb{R}_+$ &\\
       \bottomrule
        \end{tabular}
        \label{fig:placeholder}
    \end{table}
\end{enumerate}
\textbf{Zadanie domowe.} Zapisz szkic rozwiązania zadania rozwiązanego przez Ciebie przy tablicy w zeszłym tygodniu. Jeśli nie ma takiego zadania, zapisz dowolne twierdzenie, fakt lub lemat ze szkicem dowodu z któregoś z zeszłotygodniowych wykładów (wskazówka \verb|\newtheorem{twierdzenie}{Twierdzenie}| lub analogicznie \texttt{Fakt}, \texttt{Lemat}).
\end{document}
\endenumerate